\begin{foldable} % Model compression --> Knowledge distillation
    An inevitable side effect of advancing artificial intelligence (AI) that often gets overlooked is the sheer size of AI models.
    This obstacle limits their deployability in constrained environments, such as mobile or edge devices.
    Fortunately, research on model compression addresses this problem through various techniques such as model pruning, quantization, or low-rank factorization.
    Among these techniques, \textit{knowledge distillation} (KD), proposed by \citet{hinton2015distilling}, is well known for its efficiency and intuitive approach.
\end{foldable}

\begin{foldable} % Knowledge distillation
    The intuition of KD closely resembles human learning: the knowledge is transferred from an expert network (\textit{teacher}) to a novice network (\textit{student}).
    In traditional machine learning, the student learns to predict the ground-truth labels from a dataset, much like a human learner does self-studying given a collection of exercises and solutions.
    KD advantageously supplements this learning process with the teacher's knowledge, which is more informative than the ground-truth labels.
    Thus, the student can reach comparable performance to the teacher, even when it may have a lesser capacity.
\end{foldable}

\begin{foldable} % Data-free KD
    However, KD methods often require the teacher's original training data for optimal knowledge transfer.
    Vision datasets often exhibit broad semantic variation, complex environments, and substantial intra-class diversity.
    This holds true even in privacy-sensitive domains or medical domains, due to the subject's condition, environments, or equipments.
    Nonetheless, such source datasets are often protected by stringent regulations as they are tied to security, privacy, and proprietary concerns.
    In fact, personal and health data, or data as intellectual properties all qualify for protection under the Database Directive (Directive 96/9/EC) and General Data Protection Regulation (GDPR) by \citet{eu1996databasedirective, eu2016gdpr} in the EU, and under the Health Insurance Portability and Accountability Act (HIPAA) and the Data Security and Transparency Act (DSTA) by \citet{us1996hipaa, us2016dsta} in the US.
    These regulations effectively inhibit the use of traditional KD methods and pose the more challenging \textit{data-free KD} problem.
\end{foldable}

\begin{foldable} % Black-box data-free KD
    Also crucial to KD is accessibility to the teacher model, \ie, what interfaces it exposes.
    Model's accessibility vary from \textit{white-box} --- accessing the intermediate features, parameters, activations, or gradients \citep{romero2015fitnets, yim2017gift, chen2019data} to \textit{black-box} --- only the final logits or predictive probabilities \citep{wang2020neural, nguyen2022black, vo2024improving}.
    Unfortunately, the teacher's accessibility is often restricted.
    For instance, proprietary vision models, \eg, Vision AI \citep{google-visionai}, Rekognition \citep{aws-rekognition} and Azure Vision \citep{azure-vision}, return final object detection features instead of intermediate representations.
    Similarly, certain evaluation servers, \eg, ImageNet by \citet{russakovsky2015imagenet}, only return the top-1 prediction, not class probabilities.
    Retrieving deeper insights via these interfaces is usually costly or even impossible.
    In conjunction, while \textit{black-box data-free KD} (BBDFKD) is a realistic problem, its contraints invalidates most KD methods and calls for a fresh approach.
\end{foldable}

\begin{foldable} % Weakness of existing methods: Diversity
    Without training data, previous methods utilize synthetic images to tackle BBDFKD.
    In ZSDB3 \citep{wang2021zero}, synthetic images are sampled from random uniform noise and then updated as trainable parameters.
    However, they have a suboptimal initialization strategy, making some classes unforgivingly unobtainable.
    Differently, IDEAL \citep{zhang2022ideal} and DFHL-RS \citep{yuan2024data} leverage a generator network for synthetic images.
    Both attempted to generate class-balanced and realistic images, but were at risk of mode collapse caused by generator overfitting, as pointed out in \citet{yuan2024data}.
    This lack of diversity in synthetic images is detrimental to KD performance, as it limits the knowledge domain the student can learn from.
    In addition, the majority of these methods are restricted to using a single index instead of probabilities as label during KD, limiting the distillation signals.
\end{foldable}

\begin{foldable} % Our method
    \textbf{Our method.} We propose \textit{Diverse Image Priors Knowledge Distillation} (\textbf{\methodname}) to tackle the BBDFKD problem, \ie, distilling the student from a black-box teacher which only returns top-1 prediction and without access to any real data.
    We focus on \textit{diversity}, a pivotal factor to perform effective KD but was overlooked in previous methods.
    For a rigorous definition, we refer to diversity as ``\textit{how broadly the samples distribute in the data space}''.
    We induce diversity into our KD framework through three phases:
        \textit{Synthesis}: to craft a pipeline of synthetic images with diverse patterns and semantics,
        \textit{Contrast}: to optimize them to be even more distinct through contrastive learning,
        and \textit{Distillation}: to train a high-performance student from these images with soft distillation.
\end{foldable}

\begin{foldable} % Contribution
    \textbf{Contribution.} In summary, we highlight our contributions as follows:
    \begin{enumerate}
        \item We propose \textit{image priors}, synthetic images crafted with diversity in hierarchical structures, nonlinearity, and meaningful semantics.
        \item We investigate contrastive learning as a fitting strategy to improve image diversity, where synthetic images are optimized to be distinguishable to each other.
        \item We introduce a primer student to extract soft-probability as informative knowledge for KD, which is absent in previous BBDFKD methods.
        \item Our analysis demonstrates our image priors are diverse in terms of both visuals and semantics, and show that they significantly contribute to our state-of-the-art KD performance.
    \end{enumerate}
\end{foldable}
